% ===========================================================================
% Drop-in LaTeX snippets for the ICECS 2026 paper (conf branch, Exp 11).
% Numbers: results/log_conf_latency.md + results/conf_timing_results.txt.
% All latencies measured on the B-U585I-IOT02A (STM32U585, Cortex-M33 @
% 160 MHz) with the DWT cycle counter over the 1,662-frame test utterance
% p232_009, identical methodology to the baseline numbers in the draft.
% ===========================================================================

% ---------------------------------------------------------------------------
% TABLE: latency-reduction variants
% ---------------------------------------------------------------------------
\begin{table}[!t]
\caption{Per-frame latency of latency-reduction variants on the STM32U585.
Real-time factor (RTF) is against the 2.5~ms output stride. ``Lossless''
marks changes that preserve the trained network function exactly; the
others require retraining to re-establish enhancement quality.}
\label{tab:latency-variants}
\centering
\setlength{\tabcolsep}{4pt}
\begin{tabular}{lccc}
\hline
\textbf{Variant ($N{=}64$)} & \textbf{ms/frame} & \textbf{RTF} & \textbf{Lossless} \\
\hline
Baseline (as deployed)                  & 2.715          & 1.086          & ---          \\
$-$~binarization                        & 2.683          & 1.073          & no           \\
$-$~both $1\times1$ convolutions        & 2.282          & 0.913          & no           \\
ALIF$\rightarrow$PLIF readout           & 2.723          & 1.089          & no           \\
All three                               & 2.253          & 0.901          & no           \\
Decoder ConvTranspose$\rightarrow$Gemm  & 2.564          & 1.025          & \textbf{yes} \\
Decoder Gemm $+$ all three              & \textbf{2.101} & \textbf{0.840} & no           \\
\hline
$N{=}128$, decoder Gemm (lossless)      & 5.991          & 2.396          & \textbf{yes} \\
\hline
\end{tabular}
\end{table}

% ---------------------------------------------------------------------------
% METHODS paragraph (add to Sec. II, after the streaming reformulation)
% ---------------------------------------------------------------------------
% \subsection{Latency Reduction}
In the streaming formulation the decoder's transposed convolution acts on a
single feature vector per inference, so
$\mathrm{ConvTranspose1d}(N,1,K{=}80,\mathrm{stride}{=}40)$ degenerates to a
matrix--vector product: $\mathbf{y} = \mathbf{W}^{\top}\mathbf{x} + b$ with
$\mathbf{W} \in \mathbb{R}^{N\times80}$ a transposed view of the trained
kernel. We therefore re-express the decoder as a dense (Gemm) layer with
identical weights. This rewrite is mathematically lossless and removes the
zero-stuffed upsampling buffer that X-CUBE-AI otherwise materialises: peak
activation memory drops from 23.3~KB to 4.3~KB at $N{=}64$ (47~KB to 8.7~KB
at $N{=}128$), and the on-device output matches the unmodified network to
$0.00$~dB SI-SNR. In addition, we evaluate three architectural
simplifications of the separator: removing the learned-threshold
binarization, removing the two $1\times1$ convolutions (the bottleneck
projection after LayerNorm and the mask expansion before the sigmoid, valid
since $N{=}B$), and replacing the ALIF readout with a PLIF readout (scalar
instead of per-channel decay). These three change the network function and
are evaluated for latency only, without retraining.

% ---------------------------------------------------------------------------
% RESULTS paragraph (add to Sec. III, after the latency paragraph)
% ---------------------------------------------------------------------------
% \subsubsection{Reaching real time}
Table~\ref{tab:latency-variants} reports the measured per-frame latency of
the latency-reduction variants. The lossless decoder rewrite brings
$N{=}64$ from $2.715$ to $2.564$~ms per frame at unchanged quality
($16.70$~dB test SI-SNR; on-device output identical to within $0.00$~dB).
Of the three architectural simplifications, removing the two $1\times1$
convolutions is the most effective and alone crosses the real-time
threshold ($2.282$~ms, RTF~$0.913$); removing the binarization saves a
further $1.2\%$, while exchanging the ALIF readout for a PLIF readout is
latency-neutral. With all changes combined the network runs at
$2.101$~ms per frame (RTF~$0.840$), $16\%$ under the real-time budget on a
$160$~MHz Cortex-M33. Notably, the static compiler report does not predict
these gains: the decoder accounts for $95\%$ of reported multiply-accumulate
operations but only ${\sim}24$\,K cycles ($\sim$$6\%$) of measured runtime,
whereas the two $1\times1$ convolutions, billed at under $2\%$ of MACC,
cost $16\%$ of cycles. On-device measurement, rather than operation
counting, is therefore essential when optimising small streaming graphs for
microcontrollers. The simplifications that alter the network function
require retraining to re-establish enhancement quality, which we leave to
future work; the decoder rewrite can be adopted unconditionally.
